\subsection{The unforced scale-profile estimate}
\label{unforced:profile-section}

An edge from \(n\) to \(m<n\) in \([0,N]\) is admissible when
\(2n-m\le N\), and its cost is
$$
 c_g(m,n)=g(m)-\min_{[m,n]}g.
$$
No intermediate endpoint is prescribed in this section. This distinction is
essential: the following estimate can be used in a distance argument only after
the analytic transfer has been justified for such unforced chains.

\begin{lemma}[Unforced profile bound]\label{unforced:profile}
Let \(0<a\le b\le1\), and let \(g:[0,N]\to\mathbb R\) be 1-Lipschitz
with \(ax\le g(x)\le bx\). For every \(0\le n_0<N\), an admissible
decreasing chain from \(n_0\) to zero exists with at most
$$
 K\le C\left(1+\log\frac{N}{N-n_0}\right)
$$
edges and total cost at most \(NU(a,b)\), where
\begin{equation}\label{unforced:cost}
 U(a,b)=\begin{cases}
 \dfrac{(b-a)(1-a)}{2(2b-a)},&b\le1/2,\\[7pt]
 \dfrac{b-a}{1+2b},&b\ge1/2.
 \end{cases}
\end{equation}
The formulas agree at \(b=1/2\). On a grid of mesh \(T\) containing
\(0,n_0,N\), the endpoints can be chosen on the grid with additional
\(O(KT)\) cost. If the barriers hold with uniform additive error
\(\epsilon N\), the additional cost is \(O(K\epsilon N+KT)\).
\end{lemma}

\begin{proof}
The case \(b\ge1/2\) is the bounded-length unforced chain estimate of
Lemma~\ref{fprof:chain}, followed by the grid limit described below. We prove the additional case \(b\le1/2\),
normalizing \(N=1\). Set
$$
 q=\frac{b}{2b-a},\qquad r=2q-1=\frac{a}{2b-a}.
$$
Then \(1/2<q\le1\) and \(0<r\le1\).

First observe that every endpoint \(n\le q\) admits free continuation
to zero. For \(1/2<n\le q\),
\begin{equation}\label{unforced:zero-zone}
 g(2n-1)\le b(2n-1)\le an\le g(n),
\end{equation}
because \((2b-a)n\le b\). Choosing the leftmost minimum of \(g\)
on \([2n-1,n]\) gives an admissible zero-cost edge with endpoint less
than \(n\). When \(n\le1/2\), the edge directly to zero is admissible
and costs zero since \(g(0)=0\) and \(g\ge0\). On a finite grid the
procedure terminates. Consecutive zero-cost intervals can be merged whenever
their union is admissible, and their union remains zero-cost. Every two
consecutive unmergeable intervals more than double the complementary depth,
so the retained zero-cost chain has the required logarithmic length.
If \(a=b\), then \(q=1\), and this argument already proves the result
with zero cost. It also handles any starting point \(n_0\le q\).

For the remaining initial portion, define on \([r,1]\)
$$
 H(z)=\frac z2+\left(b-\frac12\right)r.
$$
This majorizes \(bz\), since
$$
 H(z)-bz=\left(\frac12-b\right)(z-r)\ge0.
$$
For \(x\ge q\), put
$$
 \Lambda(x)=H(2x-1)-g(x)
 =x-\frac12+\left(b-\frac12\right)r-g(x).
$$
The function \(\Lambda\) is nondecreasing, and
$$
 \Lambda(q)=H(r)-g(q)=br-g(q)=aq-g(q)\le0.
$$
Almost everywhere,
$$
 (-g')_+\le\frac{1-g'}2=\frac{\Lambda'}2.
$$
Use maximal admissible downward jumps through \(\{\Lambda>0\}\),
shortening the last jump if necessary to reach its lower boundary. Their
costs sum to at most \(\Lambda(n_0)_+/2\), because each edge cost is
bounded by the negative variation of \(g\) over its interval. This is at
most \(\Lambda(1)/2\).

For a subsequent endpoint \(n\ge q\) with \(\Lambda(n)\le0\), the
majorant is valid at \(2n-1\ge r\), and gives
$$
 g(2n-1)\le H(2n-1)\le g(n).
$$
Choose a leftmost minimum on \([2n-1,n]\) to take a zero-cost edge.
Continue until the endpoint becomes at most \(q\); it need not equal
\(q\). Below \(q\), use \eqref{unforced:zero-zone} to continue for free.
The maximal jumps in the positive region double complementary depth except
at the last jump. Compressing the zero-cost intervals as above gives the
claimed bound on the total number of edges.

Finally, \(H(1)\ge g(1)\ge a\), and therefore
$$
 \frac{\Lambda(1)}2
 \le\frac{1/2+(b-1/2)r-a}{2}
 =\frac{(b-a)(1-a)}{2(2b-a)}.
$$
This proves the cost bound.

For precision about arbitrary continuous profiles, first run the construction
on piecewise affine interpolants on grids with mesh tending to zero.
The interpolants retain both barriers and the Lipschitz constant. If
\(n_0\) is not a grid point, first jump to the nearest grid point below
it; for sufficiently small mesh this is admissible and costs at most the
mesh. Rounding
the potential sign change and the threshold \(q\) requires at most a
bounded number of grid crossings, each costing \(O(T)\). For example,
one stops the positive-potential construction at the first grid point above
its lower boundary, then takes the adjacent grid edge across that boundary.
Below \(q\), the barrier argument applies directly. The edge-count bound
is independent of the mesh. Pad endpoint lists by repeated endpoints and
extract a convergent subsequence. Interval minima vary continuously with
the profile and endpoints, and admissibility is closed. Removing repeated
endpoints gives the exact continuous statement. This also proves the
asserted fixed-mesh error bound.

For approximate barriers, clamp the profile between \(ax\) and \(bx\)
and interpolate its grid values. This preserves the 1-Lipschitz property
and changes the profile uniformly by \(O(\epsilon N+T)\). Every edge
cost changes by at most twice this amount. The resulting total error is
\(O(K\epsilon N+KT)\), as claimed.
\end{proof}

The improvement over the global slope-\(1/2\) majorant is quantified by
$$
 \frac{1/2-a}{2}-U(a,b)
 =\frac{a(1/2-b)}{2(2b-a)}\ge0\qquad(b\le1/2).
$$
The gain is strict below \(b=1/2\). This estimate is not asserted to be
optimal in the small-slope range.
